\documentclass[a4paper,11pt]{article}
\usepackage{needspace}
\usepackage[utf8]{inputenc}
\usepackage[T1]{fontenc}
\usepackage[italian]{babel}
\usepackage{lmodern} 
\usepackage{geometry}
\geometry{margin=2.5cm}
\usepackage{hyperref}
\usepackage{booktabs}
\usepackage{longtable}
\usepackage{array}
\usepackage{ragged2e}
\usepackage[table]{xcolor}
% list setup
\usepackage{enumitem}
\setlist{nosep}
\setlist[enumerate,1]{nosep, leftmargin=2.2em, labelwidth=1.8em, align=right}
\setlist[itemize,1]{nosep, leftmargin=2.2em, labelwidth=1.8em, align=right}
\setlist[itemize,2]{nosep, leftmargin=1.2em, labelwidth=1.8em, align=right}
% left rule
\usepackage[skins]{tcolorbox}

\newtcolorbox{ucbox}{
  enhanced, frame hidden, interior hidden,
  boxrule=0pt, sharp corners,
  borderline west={1.5pt}{0pt}{gray!60},
  left=8pt, right=0pt, top=2pt, bottom=2pt,
  before skip=10pt, after skip=12pt,
  before upper={\noindent},
  title={}
}
% Column types shared by all tables
\newcolumntype{P}[1]{>{\RaggedRight\arraybackslash}p{#1}}
\newcolumntype{H}[1]{>{\bfseries\RaggedRight\arraybackslash}p{#1}} % bold first column
\title{}
\author{}
\date{\today}

\begin{document}

\maketitle
\section*{Vista generale e Glossario}
Il sistema è un'applicazione desktop dotata di un'interfaccia grafica che non prevede accesso remoto o integrazioni con sistemi esterni. L'unico ad interfacciarsi con l'applicazione è l'Operatore che si occupa di inserire sull'applicazione le informazioni relative ai docenti, le aule e le prenotazioni.

{\renewcommand{\arraystretch}{1.35}

\begin{longtable}{@{}H{0.24\linewidth} P{0.70\linewidth}@{}}
\toprule
\textbf{Termine} & \textbf{Definizione} \\
\midrule
\endfirsthead

\toprule
\textbf{Termine} & \textbf{Definizione} \\
\midrule
\endhead

\midrule
\multicolumn{2}{r}{\footnotesize\itshape continua alla pagina successiva} \\
\endfoot

\bottomrule
\endlastfoot

\rowcolor{gray!10}Archivio & Insieme persistente di aule, docenti e prenotazioni. \\
\rowcolor{white}Aula & Luogo fisico soggetto alla prenotazione. \\
\rowcolor{gray!10}Aula Libera & Aula non soggetta a prenotazione \\
\rowcolor{white}Campo o Criterio non valido & Il dato inserito non è logicamente sensato \\
\rowcolor{gray!10}Conflitto & Sovrapposizione temporale con un'altra prenotazione della stessa aula o dello stesso docente. \\
\rowcolor{white}Docente & Colui che richiede la prenotazione. Identificato dall'insieme di dati: email univoca, nome e cognome. \\
\rowcolor{gray!10}Fascia oraria & Intervallo di tempo compreso tra una data e ora di inizio e una data e ora di fine. \\
\rowcolor{white}Operatore & Utente che gestisce aule, docenti e prenotazioni attraverso l'interfaccia. \\
\rowcolor{gray!10}Prenotazione & Occupazione di un aula per un tempo definito. \\
\rowcolor{white}Prenotazione ricorrente & Prenotazione che si ripete ad intervalli settimanali fino a una data di termine. \\
\end{longtable}
}

\newpage
\section*{Business Flow}

{\renewcommand{\arraystretch}{1.35}

\newlength{\tw}
\setlength{\tw}{\dimexpr\linewidth-8\tabcolsep\relax}

\begin{longtable}{@{}H{0.09\tw} P{0.17\tw} P{0.36\tw} P{0.23\tw} P{0.15\tw}@{}}
\toprule
\textbf{ID} & \textbf{Business Flow} & \textbf{Descrizione} & \textbf{Casi d'uso coinvolti} & \textbf{Priorità} \\
\midrule
\endfirsthead

\toprule
\textbf{ID} & \textbf{Business Flow} & \textbf{Descrizione} & \textbf{Casi d'uso coinvolti} & \textbf{Priorità} \\
\midrule
\endhead

\midrule
\multicolumn{5}{r}{\footnotesize\itshape continua alla pagina successiva} \\
\endfoot

\bottomrule
\endlastfoot

\rowcolor{gray!10}BF-01 & Gestione delle aule & L'Operatore consulta o cerca un'aula e può inserirla, modificarla o cancellarla. &\mbox{UC01}, \mbox{UC02}, \mbox{UC03}, \mbox{UC04}, \mbox{UC05} & MustHave \\
\rowcolor{white}BF-02 & Gestione dei docenti & L'Operatore consulta o cerca un docente e può inserirlo, modificarlo o cancellarlo &\mbox{UC06}, \mbox{UC07}, \mbox{UC08}, \mbox{UC09}, \mbox{UC10} & MustHave \\
\rowcolor{gray!10}BF-03 & Creazione di una prenotazione & L'Operatore può cercare le aule disponibili e inserisce una prenotazione singola o ricorrente; il sistema verifica le sovrapposizioni. &\mbox{UC11}, \mbox{UC12}, \mbox{UC13} & MustHave \\
\rowcolor{white}BF-04 & Consultazione e cancellazione delle prenotazioni & L'Operatore visualizza e filtra le prenotazioni, poi può cancellare una singola occorrenza o le prenotazioni future di una serie. &\mbox{UC14}, \mbox{UC15}, \mbox{UC20} & MustHave \\
\rowcolor{gray!10}BF-05 & Modifica di una prenotazione & L'Operatore seleziona una prenotazione singola o ricorrente e ne modifica le informazioni. &\mbox{UC14}, \mbox{UC16}, \mbox{UC18} & ShouldHave \\
\rowcolor{white}BF-06 & Conversione della ricorrenza & L'Operatore modifica una prenotazione singola rendendola ricorrente oppure converte una prenotazione ricorrente in singola. &\mbox{UC14}, \mbox{UC17}, \mbox{UC19} & NiceToHave \\
\end{longtable}
}

\newpage
\section*{Casi d'uso (UC)}
\subsection*{1. Gestione aula}
\begin{samepage}
\subsubsection*{UC01 -- Inserire aula}
\begin{ucbox}

\noindent\textbf{Attore:} Operatore\par
\noindent\textbf{Precondizioni:} L'archivio è stato caricato.\par

\noindent\textbf{Flusso di eventi:}
\begin{enumerate}
\item L'Operatore sceglie 'nuova aula'. \item Il sistema mostra il form di creazione. \item L'Operatore inserisce il codice univoco, nome, edificio, piano e capienza. \item L'Operatore invia il form. \item Il sistema convalida i dati. \item Il sistema crea l'aula. \item Il sistema salva l'archivio. \item Il sistema mostra la lista aggiornata. \end{enumerate}
\noindent\textbf{Flussi alternativi:}
\begin{itemize}
\item[4a.] L'Operatore annulla l'operazione
\begin{itemize}
\item Nessuna modifica effettuata all'archivio.
\item L'esecuzione riprende dal passo 1
\end{itemize}
\item[5a.] Il codice è già stato utilizzato
\begin{itemize}
\item Il sistema segnala l'errore e non salva l'operazione.
\item L'esecuzione riprende dal passo 2
\end{itemize}
\item[5b.] Uno dei campi è mancante o non valido
\begin{itemize}
\item Il sistema segnala l'errore e non salva l'operazione.
\item L'esecuzione riprende dal passo 2
\end{itemize}
\end{itemize}
\noindent\textbf{Postcondizioni:} Esiste una nuova aula\par
\end{ucbox}
\end{samepage}
\begin{samepage}
\subsubsection*{UC02 -- Modifica aula}
\begin{ucbox}

\noindent\textbf{Attore:} Operatore\par
\noindent\textbf{Precondizioni:} L'archivio è stato caricato. Esiste un'aula.\par

\noindent\textbf{Flusso di eventi:}
\begin{enumerate}
\item L'Operatore seleziona un'aula dalla lista \item L'Operatore seleziona l'opzione 'Modifica' \item Il sistema mostra i dati correnti \item L'Operatore modifica uno o più campi. \item L'Operatore conferma. \item Il sistema convalida i nuovi campi. \item Il sistema aggiorna l'archivio. \item Il sistema mostra la lista aggiornata. \end{enumerate}
\noindent\textbf{Flussi alternativi:}
\begin{itemize}
\item[5a.] L'Operatore annulla l'operazione
\begin{itemize}
\item Nessuna modifica effettuata all'archivio
\item L'esecuzione riprende dal passo 1
\end{itemize}
\item[6a.] Uno dei campi è mancante o non valido
\begin{itemize}
\item Il sistema segnala l'errore e non salva l'operazione
\item L'esecuzione riprende dal passo 4
\end{itemize}
\item[6b.] Tutti i campi rimangono immutati
\begin{itemize}
\item Il sistema non aggiorna i dati.
\item L'esecuzione riprende dal passo 4
\end{itemize}
\end{itemize}
\noindent\textbf{Postcondizioni:} I dati dell'aula sono aggiornati\par
\end{ucbox}
\end{samepage}
\begin{samepage}
\subsubsection*{UC03 -- Cancellare un'aula}
\begin{ucbox}

\noindent\textbf{Attore:} Operatore\par
\noindent\textbf{Precondizioni:} L'archivio è stato caricato. Esiste almeno un'aula.\par

\noindent\textbf{Flusso di eventi:}
\begin{enumerate}
\item L'Operatore seleziona un'aula \item L'Operatore seleziona l'opzione 'Cancella' \item Il sistema chiede conferma \item L'Operatore conferma \item Il sistema verifica l'esistenza di prenotazioni associate. Se non ne trova prosegue. \item Il sistema elimina l'aula \item Il sistema salva \item Il sistema mostra la lista aggiornata \end{enumerate}
\noindent\textbf{Flussi alternativi:}
\begin{itemize}
\item[4a.] L'Operatore annulla l'operazione
\begin{itemize}
\item Nessuna modifica effettuata all'archivio
\item L'esecuzione riprende dal passo 1
\end{itemize}
\item[5a.] Il sistema trova prenotazioni associate
\begin{itemize}
\item Il sistema visualizza le prenotazioni.
\item Il sistema chiede conferma.
\item L'Operatore conferma
\item Il sistema elimina le prenotazioni associate all'aula.
\item L'esecuzione riprende dal passo 6
\end{itemize}
\item[5b.] Il sistema trova prenotazioni associate
\begin{itemize}
\item Il sistema visualizza le prenotazioni.
\item Il sistema chiede conferma.
\item L'Operatore annulla
\item L'esecuzione riprende dal passo 1
\end{itemize}
\end{itemize}
\noindent\textbf{Postcondizioni:} L'aula non è più presente nell'archivio\par
\end{ucbox}
\end{samepage}
\begin{samepage}
\subsubsection*{UC04 -- Visualizzare la lista delle aule}
\begin{ucbox}

\noindent\textbf{Attore:} Operatore\par
\noindent\textbf{Precondizioni:} L'archivio è stato caricato.\par

\noindent\textbf{Flusso di eventi:}
\begin{enumerate}
\item L'Operatore apre la sezione Aule \item Il sistema mostra tutte le aule con codice, nome, edificio, piano e capienza \end{enumerate}
\noindent\textbf{Flussi alternativi:}
\begin{itemize}
\item[2a.] Nessuna aula è presente nell'archivio
\begin{itemize}
\item Il sistema mostra un messaggio di notifica.
\end{itemize}
\end{itemize}
\noindent\textbf{Postcondizioni:} Nessuna modifica ai dati effettuata. Il sitema mostra la lista delle aule.\par
\end{ucbox}
\end{samepage}
\begin{samepage}
\subsubsection*{UC05 -- Cercare un'aula}
\begin{ucbox}

\noindent\textbf{Attore:} Operatore\par
\noindent\textbf{Precondizioni:} L'archivio è stato caricato.\par

\noindent\textbf{Flusso di eventi:}
\begin{enumerate}
\item L'Operatore inserisce uno o più criteri di ricerca \item Il sistema mostra le aule che soddisfano tutti i criteri inseriti \end{enumerate}
\noindent\textbf{Flussi alternativi:}
\begin{itemize}
\item[1a.] Nessun criterio
\begin{itemize}
\item L'esecuzione riprende da UC04
\end{itemize}
\item[1b.] Criteri non validi
\begin{itemize}
\item Messaggio di errore che segnala il criterio non valido
\end{itemize}
\item[2a.] Nessun risultato
\begin{itemize}
\item Messaggio di errore 'nessuna aula trovata'
\end{itemize}
\end{itemize}
\noindent\textbf{Postcondizioni:} Nessuna modifica ai dati effettuata. Il sistema mostra l'elenco delle aule filtrato secondo i criteri di ricerca.\par
\end{ucbox}
\end{samepage}
\subsection*{2. Gestione docente}
\begin{samepage}
\subsubsection*{UC06 -- Inserire un docente}
\begin{ucbox}

\noindent\textbf{Attore:} Operatore\par
\noindent\textbf{Precondizioni:} L'archivio è stato caricato.\par

\noindent\textbf{Flusso di eventi:}
\begin{enumerate}
\item L'Operatore sceglie 'nuovo docente'. \item Il sistema mostra il form di creazione. \item L'Operatore inserisce l'email istituzionale univoca, nome e cognome. \item L'Operatore invia il form. \item Il sistema convalida i dati. \item Il sistema crea il docente. \item Il sistema salva l'archivio. \item Il sistema mostra la lista aggiornata. \end{enumerate}
\noindent\textbf{Flussi alternativi:}
\begin{itemize}
\item[4a.] L'Operatore annulla l'operazione
\begin{itemize}
\item Nessuna modifica effettuata all'archivio.
\item L'esecuzione riprende dal passo 1
\end{itemize}
\item[5a.] L'email è già stata utilizzata
\begin{itemize}
\item Il sistema segnala l'errore e non salva l'operazione.
\item L'esecuzione riprende dal passo 3
\end{itemize}
\item[5b.] Uno dei campi è mancante o non valido
\begin{itemize}
\item Il sistema segnala l'errore e non salva l'operazione.
\item L'esecuzione riprende dal passo 3
\end{itemize}
\end{itemize}
\noindent\textbf{Postcondizioni:} Esiste un nuovo docente\par
\end{ucbox}
\end{samepage}
\begin{samepage}
\subsubsection*{UC07 -- Modificare un docente}
\begin{ucbox}

\noindent\textbf{Attore:} Operatore\par
\noindent\textbf{Precondizioni:} L'archivio è stato caricato. Esiste un docente.\par

\noindent\textbf{Flusso di eventi:}
\begin{enumerate}
\item L'Operatore seleziona un docente dalla lista \item L'Operatore seleziona l'opzione 'Modifica' \item Il sistema mostra i dati correnti \item L'Operatore modifica uno o più campi. \item L'Operatore conferma. \item Il sistema convalida i nuovi campi. \item Il sistema aggiorna l'archivio. \item Il sistema mostra la lista aggiornata. \end{enumerate}
\noindent\textbf{Flussi alternativi:}
\begin{itemize}
\item[5a.] L'Operatore annulla l'operazione
\begin{itemize}
\item Nessuna modifica effettuata all'archivio
\item L'esecuzione riprende dal passo 1
\end{itemize}
\item[6a.] Uno dei campi è mancante o non valido
\begin{itemize}
\item Il sistema segnala l'errore e non salva l'operazione
\item L'esecuzione riprende dal passo 4
\end{itemize}
\item[6b.] Tutti i campi rimangono immutati
\begin{itemize}
\item Il sistema non aggiorna i dati.
\item L'esecuzione riprende dal passo 4
\end{itemize}
\end{itemize}
\noindent\textbf{Postcondizioni:} I dati del docente sono aggiornati\par
\end{ucbox}
\end{samepage}
\begin{samepage}
\subsubsection*{UC08 -- Cancellare un docente}
\begin{ucbox}

\noindent\textbf{Attore:} Operatore\par
\noindent\textbf{Precondizioni:} L'archivio è stato caricato. Esiste almeno un docente.\par

\noindent\textbf{Flusso di eventi:}
\begin{enumerate}
\item L'Operatore seleziona un docente. \item L'Operatore seleziona l'opzione 'Cancella'. \item Il sistema chiede conferma \item L'Operatore conferma \item Il sistema verifica l'esistenza di prenotazioni associate. Se non ne trova prosegue. \item Il sistema elimina il docente \item Il sistema salva \item Il sistema mostra la lista aggiornata \end{enumerate}
\noindent\textbf{Flussi alternativi:}
\begin{itemize}
\item[4a.] L'Operatore annulla l'operazione
\begin{itemize}
\item Nessuna modifica effettuata all'archivio
\item L'esecuzione riprende dal passo 1
\end{itemize}
\item[5a.] Il sistema trova prenotazioni associate
\begin{itemize}
\item Il sistema visualizza le prenotazioni.
\item Il sistema chiede conferma.
\item L'Operatore conferma
\item Il sistema elimina le prenotazioni associate al docente.
\item L'esecuzione riprende dal passo 6
\end{itemize}
\item[5b.] Il sistema trova prenotazioni associate
\begin{itemize}
\item Il sistema visualizza le prenotazioni.
\item Il sistema chiede conferma.
\item L'Operatore annulla
\item L'esecuzione riprende dal passo 1
\end{itemize}
\end{itemize}
\noindent\textbf{Postcondizioni:} Il docente non è più presente nell'archivio\par
\end{ucbox}
\end{samepage}
\begin{samepage}
\subsubsection*{UC09 -- Visualizzare la lista dei docenti}
\begin{ucbox}

\noindent\textbf{Attore:} Operatore\par
\noindent\textbf{Precondizioni:} L'archivio è stato caricato.\par

\noindent\textbf{Flusso di eventi:}
\begin{enumerate}
\item L'Operatore apre la sezione Docenti \item Il sistema mostra tutti i docenti con email istituzionale, cognome e nome. \end{enumerate}
\noindent\textbf{Flussi alternativi:}
\begin{itemize}
\item[2a.] Nessun docente è presente nell'archivio
\begin{itemize}
\item Il sistema mostra un messaggio di notifica.
\end{itemize}
\end{itemize}
\noindent\textbf{Postcondizioni:} Nessuna modifica ai dati effettuata. Il sitema mostra la lista dei docenti.\par
\end{ucbox}
\end{samepage}
\begin{samepage}
\subsubsection*{UC10 -- Cercare un docente}
\begin{ucbox}

\noindent\textbf{Attore:} Operatore\par
\noindent\textbf{Precondizioni:} L'archivio è stato caricato. L'archivio è caricato.\par

\noindent\textbf{Flusso di eventi:}
\begin{enumerate}
\item L'Operatore inserisce uno o più criteri di ricerca \item Il sistema mostra i docenti che soddisfano tutti i criteri inseriti \end{enumerate}
\noindent\textbf{Flussi alternativi:}
\begin{itemize}
\item[1a.] Nessun criterio
\begin{itemize}
\item L'esecuzione riprende da UC09
\end{itemize}
\item[1b.] Criteri non validi
\begin{itemize}
\item Messaggio di errore che segnala il criterio non valido
\item L'esecuzione riprende dal passo 1
\end{itemize}
\item[2a.] Nessun risultato
\begin{itemize}
\item Messaggio di errore 'nessun docente trovato'
\end{itemize}
\end{itemize}
\noindent\textbf{Postcondizioni:} Nessuna modifica ai dati effettuata. Il sistema mostra l'elenco dei docenti filtrato secondo i criteri di ricerca.\par
\end{ucbox}
\end{samepage}
\subsection*{3. Gestione prenotazioni}
\begin{samepage}
\subsubsection*{UC11 -- Cercare aule disponibili}
\begin{ucbox}

\noindent\textbf{Attore:} Operatore\par
\noindent\textbf{Precondizioni:} L'archivio è stato caricato. Esiste almeno un'aula.\par

\noindent\textbf{Flusso di eventi:}
\begin{enumerate}
\item L'Operatore inserisce i criteri di ricerca quali data, ora di inizio e ora di fine e/o capienza. \item Il sistema calcola le aule libere secondo i criteri inseriti. \item Il sistema mostra l'elenco delle aule libere \end{enumerate}
\noindent\textbf{Flussi alternativi:}
\begin{itemize}
\item[1a.] Criteri non validi
\begin{itemize}
\item Visualizza un messaggio di errore
\item L'esecuzione riprende dal passo 1
\end{itemize}
\item[3a.] Nessuna aula libera
\begin{itemize}
\item Visualizza un messaggio di errore
\end{itemize}
\end{itemize}
\noindent\textbf{Postcondizioni:} Il sistema mostra l'elenco delle aule libere\par
\end{ucbox}
\end{samepage}
\begin{samepage}
\subsubsection*{UC12 -- Inserire una prenotazione singola}
\begin{ucbox}

\noindent\textbf{Attore:} Operatore\par
\noindent\textbf{Precondizioni:} L'archivio è stato caricato. Esistono almeno un'aula e un docente. \par

\noindent\textbf{Flusso di eventi:}
\begin{enumerate}
\item L'Operatore sceglie 'Nuova prenotazione' \item Il sistema mostra il form. \item L'Operatore compila il form inserendo data e ora di inizio e fine attività, tipo di attività e denominazione, l'aula e il docente responsabile. \item Il sistema convalida i dati e verifica le sovrapposizioni \item Se non ci sono conflitti, il sistema crea la prenotazione e salva \end{enumerate}
\noindent\textbf{Flussi alternativi:}
\begin{itemize}
\item[3a.] Campi inseriti non validi
\begin{itemize}
\item Visualizza un messaggio di errore
\item Non crea alcuna prenotazione
\item L'esecuzione riprende dal passo 2
\end{itemize}
\item[3b.] L'Operatore seleziona 'prenotazione ricorrente'
\begin{itemize}
\item L'esecuzione riprende da UC13
\end{itemize}
\item[4a.] Si verifica un conflitto
\begin{itemize}
\item La prenotazione non è salvata
\item Viene visualizzato un messaggio di errore
\item L'esecuzione riprende dal passo 2
\end{itemize}
\end{itemize}
\noindent\textbf{Postcondizioni:} La prenotazione singola è stata creata.\par
\end{ucbox}
\end{samepage}
\begin{samepage}
\subsubsection*{UC13 -- Inserire una prenotazione ricorrente}
\begin{ucbox}

\noindent\textbf{Attore:} Operatore\par
\noindent\textbf{Precondizioni:} L'archivio è stato caricato. Esistono almeno un'aula e un docente. L'Operatore seleziona 'prenotazione ricorrente' in UC12\par

\noindent\textbf{Flusso di eventi:}
\begin{enumerate}
\item L'Operatore indica la data di termine \item Il sistema genera le occorrenze: stesso giorno della settimana e stessa fascia oraria, dalla data di inizio alla data di termine inclusa. \item Il sistema verifica le sovrapposizioni per ogni occorrenza \item Se nessuna occorrenza è in conflitto il sistema crea l'intera serie e salva \end{enumerate}
\noindent\textbf{Flussi alternativi:}
\begin{itemize}
\item[1a.] Campi inseriti non validi in UC12 o UC13
\begin{itemize}
\item Visualizza un messaggio di errore
\item L'esecuzione riprende dal passo UC12.2
\end{itemize}
\item[3a.] Il sistema rileva delle sovrapposizioni
\begin{itemize}
\item Visualizza un messaggio di errore elencando le date in conflitto
\item Non effettua la creazione della serie di eventi
\item L'esecuzione riprende dal passo UC12.2
\end{itemize}
\end{itemize}
\noindent\textbf{Postcondizioni:} La prenotazione ricorrente è stata creata.\par
\end{ucbox}
\end{samepage}
\begin{samepage}
\subsubsection*{UC14 -- Visualizzare le prenotazioni}
\begin{ucbox}

\noindent\textbf{Attore:} Operatore\par
\noindent\textbf{Precondizioni:} L'archivio è stato caricato.\par

\noindent\textbf{Flusso di eventi:}
\begin{enumerate}
\item L'Operatore apre la sezione 'Prenotazioni' \item Il sistema mostra le prenotazioni ordinate per data e ora \item L'Operatore può filtrare per aula e per docente \item Il sistema aggiorna la lista \end{enumerate}
\noindent\textbf{Flussi alternativi:}
\begin{itemize}
\item[3a.] La visualizzazione filtrata non trova alcuna corrispondenza
\begin{itemize}
\item Visualizza un messaggio di errore
\end{itemize}
\end{itemize}
\noindent\textbf{Postcondizioni:} l sitema mostra la lista delle prenotazioni.\par
\end{ucbox}
\end{samepage}
\begin{samepage}
\subsubsection*{UC15 -- Cancellare la prenotazione singola}
\begin{ucbox}

\noindent\textbf{Attore:} Operatore\par
\noindent\textbf{Precondizioni:} L'archivio è stato caricato. Esiste almeno una prenotazione.\par

\noindent\textbf{Flusso di eventi:}
\begin{enumerate}
\item L'Operatore seleziona una prenotazione (\textbf{UC14})\item L'Operatore seleziona l'opzione 'Cancella' \item Il sistema chiede conferma \item L'Operatore conferma \item Il sistema elimina solo quella occorrenza e salva. \end{enumerate}
\noindent\textbf{Flussi alternativi:}
\begin{itemize}
\item[2a.] La prenotazione appartiene a una serie
\begin{itemize}
\item Il sistema mostra la scelta 'elimina prenotazione singola' o 'elimina serie'
\item L'Operatore sceglie 'elimina prenotazione singola'
\item L'esecuzione riprende dal passo 3
\end{itemize}
\item[2b.] La prenotazione appartiene a una serie
\begin{itemize}
\item Il sistema mostra la scelta 'elimina prenotazione singola' o 'elimina serie'
\item L'Operatore sceglie 'elimina serie'
\item L'esecuzione riprende dal passo UC20.4
\end{itemize}
\item[4a.] L'Operatore annulla
\begin{itemize}
\item L'esecuzione riprende dal passo 1
\end{itemize}
\end{itemize}
\noindent\textbf{Postcondizioni:} La prenotazione non è più presente; le altre occorrenze della serie restano.\par
\end{ucbox}
\end{samepage}
\begin{samepage}
\subsubsection*{UC16 -- Modificare la prenotazione singola}
\begin{ucbox}

\noindent\textbf{Attore:} Operatore\par
\noindent\textbf{Precondizioni:} L'archivio è stato caricato. Esiste almeno una prenotazione singola.\par

\noindent\textbf{Flusso di eventi:}
\begin{enumerate}
\item L'Operatore seleziona una prenotazione singola (\textbf{UC14})\item L'Operatore seleziona l'opzione 'Modifica' \item Il sistema mostra i dati correnti \item L'Operatore modifica i campi \item L'Operatore conferma \item Il sistema verifica la validità dei dati inseriti \item Il sistema aggiorna i dati \end{enumerate}
\noindent\textbf{Flussi alternativi:}
\begin{itemize}
\item[4a.] L'Operatore modifica la prenotazione da singola a ricorrente
\begin{itemize}
\item L'esecuzione riprende da UC17
\end{itemize}
\item[5a.] L'Operatore annulla
\begin{itemize}
\item L'esecuzione riprende dal passo 1
\end{itemize}
\end{itemize}
\noindent\textbf{Postcondizioni:} La prenotazione è stata modificata\par
\end{ucbox}
\end{samepage}
\begin{samepage}
\subsubsection*{UC17 -- Modificare la prenotazione da singola a ricorrente}
\begin{ucbox}

\noindent\textbf{Attore:} Operatore\par
\noindent\textbf{Precondizioni:} L'archivio è stato caricato. Esiste almeno una prenotazione singola. L'Operatore sta modificando una prenotazione singola in UC16.\par

\noindent\textbf{Flusso di eventi:}
\begin{enumerate}
\item L'Operatore seleziona 'prenotazione ricorrente'. \item L'Operatore indica la data di termine \item Il sistema genera le occorrenze: stesso giorno della settimana e stessa fascia oraria, dalla data di inizio alla data di termine inclusa. \item Il sistema verifica le sovrapposizioni per ogni occorrenza \item Se nessuna occorrenza è in conflitto il sistema crea l'intera serie e salva \end{enumerate}
\noindent\textbf{Flussi alternativi:}
\begin{itemize}
\item[1a.] Campi inseriti non validi
\begin{itemize}
\item Visualizza un messaggio di errore
\item L'esecuzione riprende dal passo 1
\end{itemize}
\item[1b.] L'Operatore deseleziona 'prenotazione ricorrente'
\begin{itemize}
\item Nessuna modifica effettuata da UC17
\item L'esecuzione riprende da UC16
\end{itemize}
\item[2a.] Campi non validi
\begin{itemize}
\item Visualizza un messaggio di errore
\item L'esecuzione riprende dal passo 1
\end{itemize}
\item[4a.] Il sistema rileva delle sovrapposizioni
\begin{itemize}
\item Visualizza un messaggio di errore elencando le date in conflitto
\item Non effettua la creazione della serie di eventi
\item L'esecuzione riprende dal passo 1
\end{itemize}
\end{itemize}
\noindent\textbf{Postcondizioni:} Modificata la prenotazione da singola a ricorrente.\par
\end{ucbox}
\end{samepage}
\begin{samepage}
\subsubsection*{UC18 -- Modificare la prenotazione ricorrente}
\begin{ucbox}

\noindent\textbf{Attore:} Operatore\par
\noindent\textbf{Precondizioni:} L'archivio è stato caricato. Esiste almeno una prenotazione ricorrente.\par

\noindent\textbf{Flusso di eventi:}
\begin{enumerate}
\item L'Operatore seleziona una prenotazione ricorrente (\textbf{UC14})\item L'Operatore seleziona l'opzione 'Modifica' \item Il sistema mostra i dati correnti \item L'Operatore modifica i campi \item L'Operatore conferma \item Il sistema verifica la validità dei dati inseriti \item Il sistema aggiorna i dati \end{enumerate}
\noindent\textbf{Flussi alternativi:}
\begin{itemize}
\item[4a.] L'Operatore modifica la prenotazione da ricorrente a singola
\begin{itemize}
\item L'esecuzione riprende da UC19
\end{itemize}
\item[5a.] L'Operatore annulla
\begin{itemize}
\item L'esecuzione riprende dal passo 1
\end{itemize}
\item[6a.] L'Operatore ha modificato la data di termine anticipandola
\begin{itemize}
\item Il sistema verifica se ci sono prenotazioni da rimuovere
\item Il sistema avvisa e chiede conferma
\item Se l'Operatore conferma si procede con l'eliminazione delle prenotazioni e l'esecuzione riprende dal passo 6.
\item Se l'Operatore annulla l'esecuzione riprende dal passo 3.
\end{itemize}
\item[6b.] L'Operatore ha modificato la data di termine posticipandola
\begin{itemize}
\item Il sistema verifica se è possibile effettuare le nuove prenotazioni
\item Se possibile, il sistema avvisa e chiede conferma: se l'Operatore conferma si procede con la creazione degli eventi, se annulla l'esecuzione riprende dal passo 3.
\item Se non è possibile, il sistema avvisa e l'operazione viene annullata, l'esecuzione riprende dal passo 3.
\end{itemize}
\end{itemize}
\noindent\textbf{Postcondizioni:} La prenotazione ricorrente è stata modificata\par
\end{ucbox}
\end{samepage}
\begin{samepage}
\subsubsection*{UC19 -- Modificare la prenotazione da ricorrente a singola}
\begin{ucbox}

\noindent\textbf{Attore:} Operatore\par
\noindent\textbf{Precondizioni:} L'archivio è stato caricato. Esiste almeno una prenotazione ricorrente. L'Operatore sta modificando una prenotazione ricorrente in UC18\par

\noindent\textbf{Flusso di eventi:}
\begin{enumerate}
\item L'Operatore deseleziona 'prenotazione ricorrente'. \item Il sistema verifica se esistono prenotazioni successive, in tal caso le mostra e avvisa che saranno eliminate. \item L'Operatore conferma. \item Il sistema verifica la correttezza dei dati \item Il sistema elimina le occorrenze con data non uguale alla data dell'occorrenza modificata. \item Il sistema aggiorna i dati. \end{enumerate}
\noindent\textbf{Flussi alternativi:}
\begin{itemize}
\item[1a.] Campi inseriti non validi
\begin{itemize}
\item Visualizza un messaggio di errore
\item L'esecuzione riprende dal passo 1
\end{itemize}
\item[1b.] L'Operatore seleziona 'prenotazione ricorrente'
\begin{itemize}
\item Nessuna modifica effettuata da UC19
\item L'esecuzione riprende da UC18
\end{itemize}
\item[3a.] L'Operatore annulla
\begin{itemize}
\item L'esecuzione riprende dal passo UC18.1
\end{itemize}
\end{itemize}
\noindent\textbf{Postcondizioni:} La prenotazione è stata modificata da ricorrente a singola.\par
\end{ucbox}
\end{samepage}
\begin{samepage}
\subsubsection*{UC20 -- Cancellare tutte le prenotazioni future di una serie}
\begin{ucbox}

\noindent\textbf{Attore:} Operatore\par
\noindent\textbf{Precondizioni:} L'archivio è stato caricato. Esiste una prenotazione appartenente a una serie ricorrente.\par

\noindent\textbf{Flusso di eventi:}
\begin{enumerate}
\item L'Operatore seleziona una prenotazione della serie. (\textbf{UC14})\item L'Operatore seleziona l'opzione 'Cancella'. \item L'Operatore seleziona 'elimina serie'. \item Il sistema mostra le occorrenze che saranno eliminate. \item Il sistema chiede conferma. \item L'Operatore conferma. \item Il sistema elimina le occorrenze con data successiva o uguale alla data dell'occorrenza eliminata \item Il sistema aggiorna i dati. \end{enumerate}
\noindent\textbf{Flussi alternativi:}
\begin{itemize}
\item[4a.] Non ci sono occorrenze future
\begin{itemize}
\item Visualizza un messaggio
\item L'esecuzione riprende dal passo 5
\end{itemize}
\item[6a.] L'Operatore annulla
\begin{itemize}
\item L'esecuzione riprende dal passo 1
\end{itemize}
\end{itemize}
\noindent\textbf{Postcondizioni:} Sono state eliminate le occorrenze con data successiva o uguale alla data dell'occorrenza selezionata\par
\end{ucbox}
\end{samepage}



\newpage
\section*{Requisiti funzionali}

{\renewcommand{\arraystretch}{1.35}

\begin{longtable}{@{}H{0.10\linewidth} P{0.46\linewidth} P{0.12\linewidth} P{0.20\linewidth}@{}}
\toprule
\textbf{ID} & \textbf{Requisito} & \textbf{Priorità} & \textbf{Caso d'uso} \\
\midrule
\endfirsthead

\toprule
\textbf{ID} & \textbf{Requisito} & \textbf{Priorità} & \textbf{Caso d'uso} \\
\midrule
\endhead

\midrule
\multicolumn{4}{r}{\footnotesize\itshape continua alla pagina successiva} \\
\endfoot

\bottomrule
\endlastfoot

\rowcolor{gray!10}FI-01 & Il sistema deve consentire di inserire un'aula con codice univoco, nome, edificio, piano e capienza. & MustHave &\mbox{UC01} \\
\rowcolor{white}FI-02 & Il sistema deve consentire di modificare i dati di un'aula & ShouldHave &\mbox{UC02} \\
\rowcolor{gray!10}FI-03 & Il sistema deve consentire l'eliminazione di un'aula & ShouldHave &\mbox{UC03} \\
\rowcolor{white}FI-04 & Il sistema deve visualizzare la lista delle aule & MustHave &\mbox{UC04} \\
\rowcolor{gray!10}FI-05 & Il sistema deve permettere la ricerca di aule per nome, codice o capienza minima & ShouldHave &\mbox{UC05} \\
\rowcolor{white}FI-06 & Il sistema deve consentire di inserire un docente con e-mail istituzionale univoca, nome e cognome & MustHave &\mbox{UC06} \\
\rowcolor{gray!10}FI-07 & Il sistema deve consentire di modificare nome e cognome di un docente & ShouldHave &\mbox{UC07} \\
\rowcolor{white}FI-08 & Il sistema deve consentire di cancellare un docente & NiceToHave &\mbox{UC08} \\
\rowcolor{gray!10}FI-09 & Il sistema deve visualizzare la lista dei docenti ordinata per cognome e nome & NiceToHave &\mbox{UC09} \\
\rowcolor{white}FI-10 & Il sistema deve permettere la ricerca di docenti per cognome, nome o per email & NiceToHave &\mbox{UC10} \\
\rowcolor{gray!10}FI-11 & Il sistema deve permettere di cercare le aule disponibili in una data e fascia oraria & NiceToHave &\mbox{UC11} \\
\rowcolor{white}FI-12 & Il sistema deve consentire di inserire una prenotazione singola con aula, docente, tipo e denominazione attività, data e ora di inizio e fine & MustHave &\mbox{UC12} \\
\rowcolor{gray!10}FI-13 & Il sistema deve consentire di inserire una prenotazione con ricorrenza settimanale e data di termine & MustHave &\mbox{UC13} \\
\rowcolor{white}FI-14 & Il sistema deve visualizzare le prenotazioni ordinate per data e ora & ShouldHave &\mbox{UC14} \\
\rowcolor{gray!10}FI-15 & Il sistema deve permettere di filtrare le prenotazioni per aula o per docente & NiceToHave &\mbox{UC14} \\
\rowcolor{white}FI-16 & Il sistema deve consentire di cancellare una singola prenotazione & ShouldHave &\mbox{UC15} \\
\rowcolor{gray!10}FI-17 & Il sistema deve consentire di cancellare tutte le prenotazioni future di una serie ricorrente & ShouldHave &\mbox{UC20} \\
\rowcolor{white}ED-18 & Il sistema non deve ammettere due prenotazioni sovrapposte per la stessa aula & MustHave &\mbox{UC12}, \mbox{UC13}, \mbox{UC17} \\
\rowcolor{gray!10}ED-19 & Il sistema non deve ammettere due prenotazioni sovrapposte per lo stesso docente & MustHave &\mbox{UC12}, \mbox{UC13}, \mbox{UC17} \\
\rowcolor{white}ED-20 & Il sistema deve garantire l'unicità del codice aula e dell'e-mail docente & MustHave &\mbox{UC01}, \mbox{UC06} \\
\rowcolor{gray!10}ED-21 & Il sistema deve validare la correttezza dei dati in ingresso & MustHave &\mbox{UC01}, \mbox{UC02}, \mbox{UC06}, \mbox{UC07}, \mbox{UC12}, \mbox{UC13} \\
\rowcolor{white}UI-22 & Il sistema deve chiedere conferma prima di ogni cancellazione & NiceToHave &\mbox{UC03}, \mbox{UC08}, \mbox{UC15}, \mbox{UC18}, \mbox{UC19}, \mbox{UC20} \\
\rowcolor{gray!10}UI-23 & Il sistema deve mostrare messaggi di errore comprensibili qualora il sistema necessiti di segnalare un errore & NiceToHave &\mbox{UC01}, \mbox{UC02}, \mbox{UC04}, \mbox{UC05}, \mbox{UC06}, \mbox{UC07}, \mbox{UC09}, \mbox{UC10}, \mbox{UC11}, \mbox{UC12}, \mbox{UC13}, \mbox{UC14}, \mbox{UC15}, \mbox{UC20} \\
\rowcolor{white}ED-24 & Il sistema deve caricare l'archivio da file all'avvio e salvarlo dopo ogni modifica & MustHave & \\
\rowcolor{gray!10}FI-25 & Il sistema deve consentire la modifica delle informazioni di una prenotazione & ShouldHave &\mbox{UC16}, \mbox{UC17}, \mbox{UC18}, \mbox{UC19} \\
\rowcolor{white}FI-26 & Il sistema deve gestire la conversione di una prenotazione singola in una prenotazione ricorrente & NiceToHave &\mbox{UC17} \\
\rowcolor{gray!10}FI-27 & Il sistema deve gestire la conversione di una prenotazione ricorrente in una prenotazione singola & NiceToHave &\mbox{UC19} \\
\rowcolor{white}ED-28 & Il sistema deve memorizza i dati dell'aula quali Codice identificativo, nome, edificio, piano e capienza & MustHave & \\
\rowcolor{gray!10}ED-29 & Il sistema deve memorizza i dati del docente quali Email istituzionale, nome e cognome & MustHave & \\
\rowcolor{white}ED-30 & Il sistema deve memorizza i dati delle prenotazioni quali l'aula prenotata, il docente responsabile, il tipo dell'attività, la denominazione dell'attività, la data e l'ora di inizio, la data e l'ora di fine & MustHave & \\
\rowcolor{gray!10}UI-31 & Il sistema deve essere accessibile e funzionante su dispositivi desktop & MustHave & \\
\rowcolor{white}UI-32 & L'applicazione possiede una sezione 'Home' che visualizza per il giorno e l'ora corrente il numero di aule libere e prenotate, il numero di docenti liberi e prenotati e una visualizzazione breve delle prenotazioni in corso & NiceToHave & \\
\rowcolor{gray!10}UI-33 & L'applicazione possiede una sezione 'Aule' nella quale visualizza la lista delle aule, permette l'inserimento di una nuova aula, permette di effetuare la ricerca di un'aula, e permette la modifica e la cancellazione di un'aula selezionata & MustHave & \\
\rowcolor{white}UI-34 & L'applicazione possiede una sezione 'Docenti' nella quale visualizza la lista dei docenti, permette l'inserimento di una nuovo docente, permette di effetuare la ricerca di un docente, e permette la modifica e la cancellazione di un docente selezionata & MustHave & \\
\rowcolor{gray!10}UI-35 & L'applicazione possiede una sezione 'Prenotazioni' nella quale visualizza la lista delle prenotazioni, permette l'inserimento di una nuova prenotazione, permette di effetuare la ricerca di una prenotazione, e permette la modifica, rimodulazione e la cancellazione di una prenotazione selezionata & MustHave & \\
\end{longtable}
}
\newpage
\section*{Requisiti non funzionali}

{\renewcommand{\arraystretch}{1.35}

\begin{longtable}{@{}H{0.10\linewidth} P{0.36\linewidth} P{0.16\linewidth} P{0.12\linewidth} P{0.16\linewidth}@{}}
\toprule
\textbf{ID} & \textbf{Requisito} & \textbf{Categoria} & \textbf{Priorità} & \textbf{Verifica} \\
\midrule
\endfirsthead

\toprule
\textbf{ID} & \textbf{Requisito} & \textbf{Categoria} & \textbf{Priorità} & \textbf{Verifica} \\
\midrule
\endhead

\midrule
\multicolumn{5}{r}{\footnotesize\itshape continua alla pagina successiva} \\
\endfoot

\bottomrule
\endlastfoot

\rowcolor{gray!10}NF-01 & L'applicazione deve avere un'interfaccia grafica per tutte le funzionalità & Interfaccia grafica & MustHave & Test manuale \\
\rowcolor{white}NF-02 & L'archivio (aule, docenti e prenotazioni) deve essere salvato su file & Persistenza dei dati & MustHave & - \\
\rowcolor{gray!10}NF-03 & Il sistema non deve sovrascrivere il file esistente se il caricamento fallisce & Affidabilità & MustHave & Test \\
\rowcolor{white}NF-04 & Non è necessaria alcuna preparazione tecnica per usare il programma, i messaggi di errore devono essere chiari affinché sia possibile correggere le operazioni facilmente & Usabilità & ShouldHave & - \\
\rowcolor{gray!10}NF-05 & L'applicazione deve funzionare su differenti piattaforme desktop & Portabilità & NiceToHave & - \\
\rowcolor{white}NF-06 & L'applicazione deve funzionare senza connessione di rete o server esterni & - & MustHave & - \\
\rowcolor{gray!10}NF-07 & Il codice deve essere strutturato in moduli con alta coesione e basso accoppiamento. In particolare si mira a separare interfaccia utente e logica applicativa. & Manutenibilità & ShouldHave & - \\
\rowcolor{white}NF-08 & Le interfacce dei componenti devono essere documentate con commenti di documentazioni & Manutenibilità & NiceToHave & - \\
\rowcolor{gray!10}NF-09 & La logica applicativa deve essere coperta da test automatici & Testabilità & MustHave & - \\
\end{longtable}
}
\newpage
\section*{Matrice di tracciabilità}

{\renewcommand{\arraystretch}{1.35}

\begin{longtable}{@{}H{0.10\linewidth} P{0.20\linewidth} P{0.20\linewidth} P{0.12\linewidth} P{0.12\linewidth} P{0.12\linewidth}@{}}
\toprule
\textbf{Caso d'uso} & \textbf{Requisiti funzionali} & \textbf{Requisiti non funzionali} & \textbf{Design} & \textbf{Codice} & \textbf{Test} \\
\midrule
\endfirsthead

\toprule
\textbf{Caso d'uso} & \textbf{Requisiti funzionali} & \textbf{Requisiti non funzionali} & \textbf{Design} & \textbf{Codice} & \textbf{Test} \\
\midrule
\endhead

\midrule
\multicolumn{6}{r}{\footnotesize\itshape continua alla pagina successiva} \\
\endfoot

\bottomrule
\endlastfoot

\rowcolor{gray!10}\mbox{UC01} & \mbox{ED-21}, \mbox{UI-23}, \mbox{ED-20}, \mbox{FI-01} & \mbox{NF-01}, \mbox{NF-04} & & & \\
\rowcolor{white}\mbox{UC02} & \mbox{ED-21}, \mbox{FI-02}, \mbox{UI-23} & \mbox{NF-01}, \mbox{NF-04} & & & \\
\rowcolor{gray!10}\mbox{UC03} & \mbox{UI-22}, \mbox{FI-03} & \mbox{NF-01} & & & \\
\rowcolor{white}\mbox{UC04} & \mbox{FI-04}, \mbox{UI-23} & \mbox{NF-01}, \mbox{NF-04} & & & \\
\rowcolor{gray!10}\mbox{UC05} & \mbox{FI-05}, \mbox{UI-23} & \mbox{NF-01}, \mbox{NF-04} & & & \\
\rowcolor{white}\mbox{UC06} & \mbox{FI-06}, \mbox{ED-21}, \mbox{UI-23}, \mbox{ED-20} & \mbox{NF-01}, \mbox{NF-04} & & & \\
\rowcolor{gray!10}\mbox{UC07} & \mbox{FI-07}, \mbox{ED-21}, \mbox{UI-23} & \mbox{NF-01}, \mbox{NF-04} & & & \\
\rowcolor{white}\mbox{UC08} & \mbox{UI-22}, \mbox{FI-08} & \mbox{NF-01} & & & \\
\rowcolor{gray!10}\mbox{UC09} & \mbox{FI-09}, \mbox{UI-23} & \mbox{NF-01}, \mbox{NF-04} & & & \\
\rowcolor{white}\mbox{UC10} & \mbox{FI-10}, \mbox{UI-23} & \mbox{NF-01}, \mbox{NF-04} & & & \\
\rowcolor{gray!10}\mbox{UC11} & \mbox{UI-23}, \mbox{FI-11} & \mbox{NF-01}, \mbox{NF-04} & & & \\
\rowcolor{white}\mbox{UC12} & \mbox{FI-12}, \mbox{ED-18}, \mbox{ED-21}, \mbox{UI-23}, \mbox{ED-19} & \mbox{NF-01}, \mbox{NF-04} & & & \\
\rowcolor{gray!10}\mbox{UC13} & \mbox{FI-13}, \mbox{ED-18}, \mbox{ED-21}, \mbox{UI-23}, \mbox{ED-19} & \mbox{NF-01}, \mbox{NF-04} & & & \\
\rowcolor{white}\mbox{UC14} & \mbox{FI-15}, \mbox{UI-23}, \mbox{FI-14} & \mbox{NF-01}, \mbox{NF-04} & & & \\
\rowcolor{gray!10}\mbox{UC15} & \mbox{UI-22}, \mbox{FI-16}, \mbox{UI-23} & \mbox{NF-01} & & & \\
\rowcolor{white}\mbox{UC16} & \mbox{FI-25} & \mbox{NF-01}, \mbox{NF-04} & & & \\
\rowcolor{gray!10}\mbox{UC17} & \mbox{ED-18}, \mbox{FI-25}, \mbox{ED-19}, \mbox{FI-26} & \mbox{NF-01}, \mbox{NF-04} & & & \\
\rowcolor{white}\mbox{UC18} & \mbox{UI-22}, \mbox{FI-25} & \mbox{NF-01}, \mbox{NF-04} & & & \\
\rowcolor{gray!10}\mbox{UC19} & \mbox{UI-22}, \mbox{FI-25}, \mbox{FI-27} & \mbox{NF-01}, \mbox{NF-04} & & & \\
\rowcolor{white}\mbox{UC20} & \mbox{UI-22}, \mbox{FI-17}, \mbox{UI-23} & \mbox{NF-01}, \mbox{NF-04} & & & \\
\end{longtable}
}

\end{document}